\begin{frame}{왜 $n$이 커지면 편향은 줄고 분산은 커지는가}
  \begin{itemize}
    \item $V$가 참값 $V^\pi$라면 $G_t^{(n)}$의 기대값은
      \textbf{정확히} $V^\pi(s_t)$ --- 실제 보상 합이 ``앞으로 받을
      할인 보상의 기댓값''의 단일 샘플(모든 $n$에서 성립).
    \item 학습 초반 $V \ne V^\pi$: 목표값에 추정치가 \textbf{일부}
      섞여 있는 것 = \kb{편향(bias)}의 출처.
  \end{itemize}
  \begin{columns}[T]
    \begin{column}{0.49\textwidth}
      \kb{편향 $\downarrow$}
      $n$이 커지면 목표에서 추정치 비중 $\downarrow$, 실제 보상
      비중 $\uparrow$ $\Rightarrow$ $V$가 부정확해도 목표값이 참값에
      가까워짐.
    \end{column}
    \begin{column}{0.49\textwidth}
      \kb{분산 $\uparrow$}
      실제 보상 합이 길어지면 각 스텝의 전이·보상 무작위성이 더 많이
      누적됨. MC가 ``분산 극대''인 이유와 정확히 같음.
    \end{column}
  \end{columns}
  \vskip0.6em
  {\footnotesize ML1 Week 4.1의 \textbf{편향-분산 트레이드오프}가 여기서는
  $n$이라는 하나의 다이얼 위 숫자로 드러남.}
\end{frame}
